\subsection{标量积为负的向量系}

引理 3. 设 \(q\) 是实向量空间 \(V\) 上的一个正二次型，\(B\) 是与之关联的对称双线性型。设 \(a_1, \ldots, a_n\) 是 \(V\) 中满足 \(B(a_i, a_j) \leqslant 0\) 对 \(i \neq j\) 的元素。

\begin{enumerate}
\def\labelenumi{(\roman{enumi})}
\tightlist
\item
若 \(c_{1}, \ldots, c_{n}\) 是满足 \(q\left(\sum_{i} c_{i} a_{i}\right) = 0\) 的实数，则
\end{enumerate}

\[
q \left(\sum_ {i} | c _ {i} | a _ {i}\right) = 0.
\]

\begin{enumerate}
\def\labelenumi{(\roman{enumi})}
\setcounter{enumi}{1}
\tightlist
\item
若 \(q\) 非退化，且存在 \(f\) 上的线性形式 \(V\)，使得 \(f(a_i) > 0\) 对所有 \(i\) 成立，则向量 \(a_1, \ldots, a_n\) 线性无关。
\end{enumerate}

关系 \(\mathrm{B}(a_i, a_j) \leqslant 0\) 对 \(i \neq j\) 立即推出

\[
q \left(\sum_ {i} | c _ {i} | a _ {i}\right) \leqslant q \left(\sum_ {i} c _ {i} a _ {i}\right),
\]

若 \(q\) 非退化，则关系 \(\sum_{i} c_{i} a_{i} = 0\) 从而蕴含

\[
\sum_ {i} | c _ {i} | a _ {i} = 0;
\]

于是，对于 \(f\) 上的任意线性形式 \(\mathbf{V}\)，有 \(\sum_{i}|c_i|f(a_i) = 0\)，从而 \(c_i = 0\) 对所有 \(i\) 成立，条件是还对所有 \(f(a_i) > 0\) 有 \(i\)。这证明了 (ii)。

引理 4. 设 \(\mathbf{Q} = (q_{ij})\) 是一个实的、对称的、\(n\) 阶方阵，满足：

\begin{enumerate}
\def\labelenumi{\alph{enumi})}
\item
  \(q_{ij} \leqslant 0\) for \(i \neq j\) ;
\item
不存在将 \(\{1,2,\ldots,n\}\) 划分为两个非空子集 I 和 J 的划分，使得 \((i,j)\in\mathrm{I}\times\mathrm{J}\) 蕴含 \(q_{ij}=0\)；
\item
二次型 \(q(x_{1},\ldots ,x_{n}) = \sum_{i,j}q_{ij}x_{i}x_{j}\) 定义在 \(\mathbf{R}^n\) 上且为正。
\end{enumerate}

于是：

\begin{enumerate}
\def\labelenumi{(\roman{enumi})}
\item
\(\mathbf{N}\) 是 \(q\) 的核，其维数为 0 或 1。若 \(\dim \mathbf{N} = 1\)，则 \(\mathbf{N}\) 由一个所有坐标均 \(> 0\) 的向量生成。
\item
\(\mathbf{Q}\) 的最小特征值的重数为 1，并且相应的特征向量的所有坐标均 \(>0\) 或均 \(< 0\)。
\end{enumerate}

由于 \(q\) 是正二次型，\(N\) 是 \(q\) 的迷向向量集合（关于 \(q\) 而言）（《代数》第 IX 章 §7 第 1 号命题 2 的推论）。令 \(a_1, \ldots, a_n\) 为 \(\mathbf{R}^n\) 的标准基。若 \(\sum_{i} c_i a_i \in N\)，引理 3 表明 \(\sum_{i} |c_i| a_i \in N\) 也成立，从而有 \(\sum_{i} q_{ji}|c_i| = 0\) 对所有 \(j\) 成立。令 \(I\) 为 \(i\) 满足 \(c_i \neq 0\) 时的指标集合。若 \(j \notin I\)，则 \(q_{ji}|c_i| \leqslant 0\) 对 \(i \in I\) 成立，而 \(q_{ji}|c_i| = 0\) 对 \(i \notin I\) 成立，所以 \(q_{ji} = 0\) 对 \(j \notin I\) 且 \(i \in I\) 成立。因此，假设 \(b)\) 蕴含 \(I = \varnothing\) 或 \(I = \{1, \ldots, n\}\)。于是，\(N\) 中每个非零向量的所有坐标都 \(\neq 0\)。若 \(\dim N \geqslant 2\)，则 \(N\) 与方程为 \(x_i = 0\) 的超平面的交的维数将 \(\geqslant 1\)，这与我们刚才证明的事实矛盾。前面的论证还表明，若 \(\dim N = 1\)，则 \(N\) 含有一个所有坐标均 \(>0\) 的向量。这证明了 (i)。

另一方面，我们知道 Q 的特征值为实数（《代数》第 IX 章 §7 第 3 号命题 5），且由于 q 为正的而为正。令 \(\lambda\) 为其中最小者。矩阵 \(Q' = Q - \lambda I_n\) 是退化正型 \(q'\) 的矩阵，而 \(Q'\) 的非对角元素与 Q 的相同。因此，\(Q'\) 满足本引理陈述中的条件 a)、b) 和 c)。由于 \(N'\) 是 \(q'\) 的核，且是 Q 对应于特征值 \(\lambda\) 的特征子空间，断言 (ii) 由 (i) 得出。

引理 5. 设 \(e_1, \ldots, e_n\) 是生成 \(T\) 的向量，满足：

\(a) (e_i | e_j) \leqslant 0 \text{ 当 } i \neq j;\)

\begin{enumerate}
\def\labelenumi{\alph{enumi})}
\setcounter{enumi}{1}
\tightlist
\item
不存在将 \(\{1, \ldots, n\}\) 划分为两个非空子集 I 和 J 的划分，使得 \((e_i|e_j) = 0\) 对 \(i \in I\) 且 \(j \in J\) 成立。
\end{enumerate}

于是有两种可能：

\begin{enumerate}
\def\labelenumi{\arabic{enumi})}
\item
\((e_1, \ldots, e_n)\) 是 T 的一个基；
\item
\(n = \dim T + 1\)；存在一个族 \((c_i)_{1 \leqslant i \leqslant n}\)，其元素为 \(> 0\) 的实数，并满足 \(\sum_{i} c_i e_i = 0\)，并且任何实数族 \((c_i')_{1 \leqslant i \leqslant n}\) 满足 \(\sum_{i} c_i' e_i = 0\) 时，都与 \((c_i)_{1 \leqslant i \leqslant n}\) 成比例。
\end{enumerate}

置 \(q_{ij} = (e_i|e_j)\)。于是矩阵 \(Q = (q_{ij})\) 满足引理 4 的假设：引理 4 的条件 a) 和 b) 分别就是上述条件 a) 和 b)，而 c) 得到满足，因为 \(\sum_{i,j} q_{ij} x_i x_j = \|\sum_{i} x_i e_i\|^2\)。\(q\) 在 \(\mathbf{R}^n\) 上且矩阵为 \(\mathbf{Q}\) 的二次型的核 N，是 \((c_1, \ldots, c_n) \in \mathbf{R}^n\) 中满足 \(\sum_{i} c_i e_i = 0\) 的向量所成的集合。若 \(N = \{0\}\)，则 \(e_i\) 线性相关，我们处于情形 1)。若 \(\dim N > 0\)，引理 4(i) 表明我们处于情形 2)。

引理 6. 设 \((e_1, \ldots, e_n)\) 是 T 的一个基，且 \((e_i | e_j) \leqslant 0\) 对 \(i \neq j\) 成立。

\begin{enumerate}
\def\labelenumi{(\roman{enumi})}
\item
若 \(x = \sum_{i} c_i e_i \in \mathbf{T}\) 满足 \((x|e_i) \geqslant 0\) 对所有 \(i\) 成立，则 \(c_i \geqslant 0\) 对所有 \(i\) 成立。
\item
若 \(x\) 和 \(y\) 是 \(\mathrm{T}\) 中的两个元素，满足 \((x|e_i) \geqslant 0\) 且 \((y|e_i) \geqslant 0\) 对所有 \(i\) 成立，则 \((x|y) \geqslant 0\)。若 \((x|e_i) > 0\) 且 \((y|e_i) > 0\) 对所有 \(i\) 成立，则 \((x|y) > 0\)。
\end{enumerate}

在 (i) 的假设下，假设某个 \(c_{i} < 0\)，其中 \(i\) 为固定指标。令 \(f\) 为 \(T\) 上由 \(f(e_{i}) = 1\) 定义的线性形式，并且

\[
f \left(e _ {j}\right) = - c _ {i} / \left(\sum_ {k = 1} ^ {n} \left| c _ {k} \right|\right) \quad \text { for } j \neq i.
\]

于是，向量 \(-x, e_1, \ldots, e_n\) 满足引理 3(ii) 的假设（取 T 上的度量型作为 \(q\)）。我们得到它们线性无关，这不可能。因此 (i) 成立。此外，若 \(x = \sum_{i} c_i e_i \in \mathbf{T}\) 且 \(y \in \mathbf{T}\)，则 \((x|y) = \sum_{i} c_i (e_i | y)\)，所以 (ii) 立即由 (i) 得出。

